% !TeX root=../main.tex
\chapter{کلیات تحقیق}
\label{chap:introduction}

\section{مقدمه}
هدف این فصل معرفی مسئله پژوهش، اهمیت و کاربردهای موضوع، و جایگاه رساله در حوزهٔ شناسایی و بازسازی کور%
\footnote{بازسازی کور یعنی بازیابی ساختار ناشناختهٔ کد کانال صرفاً از روی مشاهدات دریافتی، بدون دانش قبلی از پارامترهای فرستنده.}
کدهای کانال است%
\LTRfootnote{Blind reconstruction recovers an unknown channel-code structure from received observations without prior knowledge of the transmitter parameters.}%
. در این فصل از ورود به جزئیات ریاضی الگوریتم خودداری می‌شود و تمرکز بر تبیین روشن مسئله و انگیزهٔ حل آن قرار دارد.

\section{بیان مسئله}
مسئلهٔ محوری این رساله، بازسازی کور مشترک ساختار کد توربو شامل پلی‌نوم‌های کدگذارهای RSC، درهم‌ریز داخلی و درهم‌ریز خارجی، با درنظرگرفتن خاتمهٔ ترلیس و نویز کانال در مخابرات غیرهمکارانه است. برخلاف تشخیص صرف نوع کد، هدف بازیابی پارامترهای ساختاری ناشناخته از روی مشاهدات دریافتی است.
شکل~\ref{fig:turbo-structure-sample} نمای کلی از ساختار موازی یک کدگذار توربو را نشان می‌دهد که مؤلفه‌های اصلی مسئلهٔ بازسازی در این رساله را مشخص می‌کند.

\begin{figure}[ht]
	\centering
	% برای افزودن شکل جدید: فایل تصویر را در پوشه img قرار دهید و نام آن را در دستور زیر عوض کنید.
	\includegraphics[width=0.85\textwidth]{turbo-structure-sample}
	\caption{نمای کلی ساختار موازی کدگذار توربو (نمونه برای الگوبرداری در افزودن شکل‌های بعدی)}
	\label{fig:turbo-structure-sample}
\end{figure}

\section{اهمیت و ضرورت تحقیق}
در سناریوهای غیرهمکارانه، گیرنده معمولاً به مشخصات کامل کدگذاری دسترسی ندارد. بازسازی ساختار کد توربو امکان رمزگشایی، تحلیل ترافیک و توسعهٔ سامانه‌های شنود هوشمند را فراهم می‌کند و از نظر علمی نیز به غنی‌سازی نظریهٔ شناسایی کور کد کانال می‌انجامد.

\section{کاربردهای شناسایی کور کدهای کانال}
کاربردهای عملی شامل مخابرات نظامی و اطلاعاتی، پایش طیف، مهندسی معکوس سامانه‌های مخابراتی، و بازیابی اطلاعات در شرایط فقدان مشخصات فرستنده است. کد توربو به‌عنوان یکی از کدهای رایج استاندارد، مورد مطالعاتی مرکزی این رساله است.

\section{چالش‌های شناسایی کور کدهای توربو}
فضای جست‌وجوی درهم‌ریزها به‌صورت نمایی بزرگ است، نویز قیود جبری را مخدوش می‌کند، و ابهام ساختاری ممکن است منجر به چندین بازسازی مجاز شود. علاوه بر این، خاتمهٔ ترلیس، پانچینگ و همزمان‌سازی فریم پیچیدگی مسئله را افزایش می‌دهند.

\section{اهداف تحقیق}
هدف کلی، ارائهٔ چارچوبی برای شناسایی و بازسازی کور ساختار کد توربو است. اهداف جزئی شامل صورت‌بندی ریاضی پارامترهای مجهول، طراحی الگوریتم بازسازی مشترک، تحلیل پیچیدگی و صحت، و ارزیابی تجربی عملکرد تحت نویز است.

\section{پرسش‌های اصلی تحقیق}
\begin{enumerate}
	\item چگونه می‌توان پلی‌نوم‌های RSC و درهم‌ریزهای داخلی و خارجی را به‌طور مشترک از مشاهدات نویزی بازسازی کرد؟
	\item نقش خاتمهٔ ترلیس و قیود جبری در کاهش فضای جست‌وجو چیست؟
	\item رابطهٔ سطح نویز با ابهام ساختاری و پیچیدگی محاسباتی چگونه است؟
\end{enumerate}

\section{فرضیات و محدودیت‌های تحقیق}
فرضیات پایه شامل مدل کانال AWGN با مدولاسیون BPSK، امکان تبدیل مشاهدهٔ سخت به مدل BSC معادل، و تمرکز بر ساختار موازی کد توربو با خاتمهٔ ترلیس است. محدودیت‌هایی مانند پانچینگ ناشناخته یا عدم همزمانی کامل در فصول توسعه بررسی می‌شوند.

\section{نوآوری‌های تحقیق}
\begin{enumerate}
	\item صورت‌بندی جامع مسئلهٔ شناسایی و بازسازی کور کدهای توربو؛
	\item بازسازی کور پلی‌نوم‌های RSC بدون دانش قبلی؛
	\item بازسازی همزمان درهم‌ریز داخلی و خارجی؛
	\item لحاظ‌کردن خاتمهٔ ترلیس در بازسازی؛
	\item طراحی جست‌وجوی درختی مقید؛
	\item طراحی معیار آماری برای هرس تحت نویز؛
	\item مدل‌سازی و تحلیل ابهام ساختاری؛
	\item تحلیل رابطهٔ نویز، ابهام و پیچیدگی محاسباتی.
\end{enumerate}

\section{روش تحقیق}
روش پژوهش ترکیبی از مدل‌سازی ریاضی، استخراج قیود جبری، آزمون آماری قیود، طراحی الگوریتم جست‌وجوی درختی و ارزیابی با شبیه‌سازی مونت‌کارلو است.

\section{ساختار رساله}
فصل دوم مبانی نظری و پیشینه را مرور می‌کند. فصل سوم مدل سیستم و صورت‌بندی مسئله را ارائه می‌دهد. فصل چهارم روش پیشنهادی را شرح می‌دهد. فصل پنجم به تحلیل نظری و پیچیدگی می‌پردازد. فصل ششم نتایج شبیه‌سازی را گزارش می‌کند. فصل هفتم توسعه‌ها و مطالعات تکمیلی را بررسی می‌کند و فصل هشتم جمع‌بندی و پیشنهادهای آینده را ارائه می‌دهد.
